\section{Decisiones de normalización, embebido, referencia y desnormalización}

La información que se identifica, relaciona, autoriza o filtra se mantiene
tipada y normalizada. Categorías, perfiles, clases documentales y modelos de
embedding se separan de las entidades que los usan. Las relaciones N:M usan
tablas intermedias. Así no se repiten nombres ni reglas, y las claves foráneas
rechazan referencias inválidas.

La desnormalización se reserva para hechos históricos. Cada línea de pedido
guarda el precio y el descuento aplicados, aunque cambie la condición
comercial. Un resultado estructurado conserva parámetros, contenido, instante y
hash; la evidencia documental fija el embedding usado. Estas instantáneas no
reemplazan a las entidades vigentes: guardan el estado necesario para explicar
una operación o respuesta pasada.

JSONB se usa en cuatro lugares de estructura variable acotada:
\path{version_documental.metadatos},
\path{resultado_estructurado.parametros},
\path{resultado_estructurado.contenido} y
\path{evento_auditoria.detalles}. No se crearon índices GIN porque los
nueve casos verificables no filtran por claves internas de esos objetos. Las
columnas que intervienen en integridad, autorización o joins permanecen fuera
de JSONB.

Los archivos Markdown no se guardan en la base. Cada versión conserva una ruta
relativa, nombre, tipo MIME, tamaño y SHA-256. Así PostgreSQL controla la
identidad del archivo sin guardar binarios en la tabla de versiones.
